\color{unidarkgreen}\section[Транспортирование и хранение]{Транспортирование и хранение}\color{black}

\par
Транспортирование и хранение изделия осуществляется в транспортной таре, которая предохраняет его от повреждения.
Транспортная тара не должна вскрываться до прибытия на место монтажа.

\par
Условия транспортирования, хранения шкафа и допустимые сроки сохраняемости в упаковке до ввода в эксплуатацию приведены в таблице \ref{desc:tbl1}.

\par
Транспортирование упакованного шкафа может проводиться любым видом закрытого транспорта.
При этом транспортная тара шкафа должна быть закреплена неподвижно.

\par
Погрузка, крепление и перевозка шкафа в транспортных средствах должны осуществляться в соответствии с действующими правилами перевозок грузов на соответствующих видах транспорта, причем погрузка, крепление и перевозка шкафа железнодорожным транспортом должна проводиться в соответствии с «Техническими условиями погрузки и крепления грузов» и «Правилами перевозок грузов», утвержденными Министерством путей сообщения.

\setlength{\extrarowheight}{0.06cm}
\setlength{\tabcolsep}{1.5pt}
\setlength{\LTleft}{0pt}
\setlength{\LTright}{0pt}
\small
\begin{longtable}{
  |>{\centering\arraybackslash}m{4cm}
  |>{\centering\arraybackslash}m{3.3cm}
  |>{\centering\arraybackslash}m{3.3cm}
  |>{\centering\arraybackslash}m{3.6cm}
  |>{\centering\arraybackslash}m{3cm}|
}
\caption{Условия транспортирования и хранения\hfill\vspace{-0.5\baselineskip}}\label{desc:tbl1}\\ 
\hline
\rowcolor{gray!30}
% Первая строка заголовка
Вид поставок &
\multicolumn{2}{>{\centering\arraybackslash}m{6.72cm}|}{\centering Обозначение условий транспортирования в части воздействия} &
Обозначение условий хранения в части воздействия климатических факторов по ГОСТ 15150-69 &
Допустимые сроки сохраняемости в упаковке поставщика, годы \\
% Вторая строка заголовка (подзаголовки)
\hhline{>{\arrayrulecolor{gray!30}}-|>{\arrayrulecolor{black}}-|>{\arrayrulecolor{black}}-|>{\arrayrulecolor{gray!30}}-|>{\arrayrulecolor{gray!30}}-|}
\rowcolor{gray!30}
& Механических факторов по ГОСТ 23216-78 & Климатических факторов по ГОСТ 15150-69 & & \\
\arrayrulecolor{black}
\hline
\endfirsthead

% Повтор для следующих страниц
%\caption*{\emph{Продолжение таблицы \ref{desc:tbl1}}} \\
\hline
\rowcolor{gray!30}
% Первая строка заголовка
Вид поставок &
\multicolumn{2}{>{\centering\arraybackslash}m{6.63cm}|}{\centering Обозначение условий транспортирования в части воздействия} &
Обозначение условий хранения в части воздействия климатических факторов по ГОСТ 15150-69 &
Допустимые сроки сохраняемости в упаковке поставщика, годы \\
% Вторая строка заголовка (подзаголовки)
\hhline{>{\arrayrulecolor{gray!30}}-|>{\arrayrulecolor{black}}-|>{\arrayrulecolor{black}}-|>{\arrayrulecolor{gray!30}}-|>{\arrayrulecolor{gray!30}}-|}
\rowcolor{gray!30}
& Механических факторов по ГОСТ 23216-78 & Климатических факторов по ГОСТ 15150-69 & & \\
\arrayrulecolor{black}
\hline
\endhead

\hline
\endfoot
\endlastfoot

% Тело таблицы
\raggedright{Внутри страны (кроме районов Крайнего Севера и приравненных к ним местностей по ГОСТ 15846-2002)} & С & 5(ОЖ4) & 2(С) & 3 \\ \hline
\raggedright{Внутри страны в районы Крайнего Севера и приравненные к ним местности по ГОСТ 15846-2002} & С & 5(ОЖ4) & 2(С) & 3 \\ \hline
\raggedright{Экспортные районы с умеренным климатом} & С & 5(ОЖ4) & 2(С) & 3 \\ \hline

\end{longtable}
\setlength{\tabcolsep}{6pt}

\noindent Примечания: \\
1. Нижнее значение температуры окружающего воздуха при транспортировании и хранении определяется комплектующей элементной базой и материалами, применяемыми в шкафах, и составляет минус 25 \textdegree C. \\
2. Требования по условиям хранения распространяются на склады изготовителя и потребителя продукции. \\
3. Сроки транспортирования входят в общий срок сохраняемости изделий. Сроки транспортирования и промежуточного хранения при перегрузках не должны превышать 3 мес — для условий С. \\
4. Морская перевозка изделий допускается только для отдельного вида упаковки в трюмах судов или в контейнерах. Перевозка изделий на открытой палубе судов не допускается. \\

\par
Техническое обслуживание изделий в период хранения должно включать внешний осмотр транспортной тары, контроль целостности упаковки.